\section{The railway circuit}
\label{sec:railway}

The computing model behind every universal automaton in this line of work is the railway circuit of Stewart
\cite{stewart1994}, as adapted by Margenstern \cite{margenstern2008vol2}. We recall it here in the form we
need, which is purely combinatorial. A railway is a network of tracks on which one locomotive runs. The
locomotive is never created or destroyed and never stops. The whole content of a computation is held in the
settings of the switches the locomotive flips as it rolls.

\subsection{Tracks and the locomotive}

A track is a path. The locomotive occupies two consecutive positions on it, a leading position we call the
head and the position just behind it we call the trailing cell. At each step the pair advances by one position
along the track, the head moving to the next position and the trailing cell following into the position the
head has left. The pair is what gives the locomotive a direction. A head alone could not tell forward from
backward, but a head with a trailing cell behind it can, because the way back is occupied and the way forward
is open.

\subsection{The three switches}

A switch is a place where three tracks meet. One of the three is distinguished and called the trunk. The other
two are the branches, named $a$ and $b$. A switch is an oriented object. Entering it from the trunk is an
active crossing, in which the switch sends the locomotive on to one branch, the one it currently selects.
Entering it from a branch is a passive crossing, in which the locomotive simply proceeds to the trunk. The
three kinds of switch differ only in how the selected branch is updated.

\begin{itemize}
  \item A \emph{fixed} switch never changes its selection. Entered from the trunk it always sends the
  locomotive to the same branch. It is a permanent fork.
  \item A \emph{flip-flop} switch changes its selection on every active crossing. Entered from the trunk it
  sends the locomotive to the selected branch and then selects the other branch for next time. The passive
  crossing leaves it unchanged. The flip-flop is the only one-way switch, used only in the active direction.
  \item A \emph{memory} switch records the branch through which it was last entered passively. A passive
  crossing from branch $a$ sets its selection to $a$, a passive crossing from branch $b$ sets it to $b$. An
  active crossing then sends the locomotive to the branch the switch last remembered.
\end{itemize}

A plain crossing, where two tracks cross without interacting, is also allowed. In the plane a crossing is a
genuine four-way piece. In three-dimensional space the third dimension lets one track pass over the other, so
a crossing becomes a bridge and no special piece is needed.

\subsection{From switches to a register machine}

A register holds a non-negative integer. A register unit is built from a flip-flop paired with a memory
switch, and a chain of such units holds the value in unary, as the number of units currently set. The
locomotive performs the three register operations by entering the chain through the appropriate track. To
increment, it runs in, sets the first free unit, and returns. To decrement, it runs in, clears the last set
unit, and returns. When it is asked to decrement a register that is already zero, it finds no set unit to
clear and comes back along a different track, which is how the machine learns that the register was zero. This
zero test is the one piece of control that a counter machine needs.

A sequencer wires the instructions of a program together. Each instruction of a register machine is a small
sub-circuit that drives the locomotive to the register it acts on, performs the operation, reads the zero test
where needed, and routes the locomotive to the next instruction. Loops and conditional jumps are just track
that returns to an earlier instruction. The circuit is infinite, since the registers are unbounded, but it is
ultimately periodic, a fixed pattern repeated outward, which is all that weak universality requires.

\subsection{Why the railway suits hyperbolic space}

Two facts make the railway the natural choice. First, there is no synchronisation to maintain. A single
locomotive is the entire moving part, so the exponential stretching of distances in hyperbolic space, which
defeats any scheme that must time the meeting of separate signals, costs nothing here. The locomotive simply
takes longer to reach the far parts of the circuit, and that is harmless. Second, the railway needs only local
routing decisions, made one switch at a time, and a hyperbolic tiling has room in abundance to lay out the
tracks and to give crossings the space they need.

The reduction is standard and we do not reprove it. A register machine with two counters is universal
\cite{minsky1967}. The railway with the three switches above simulates any register machine
\cite{stewart1994, margenstern2008vol2}. We record the consequence we will use.

\begin{proposition}[Railway universality]
\label{prop:railway}
Any device that faithfully realizes a track on which a head-and-trailing-cell locomotive advances in one
direction, together with the fixed, flip-flop, and memory switches and a crossing, simulates an arbitrary
register machine, and is therefore universal.
\end{proposition}

The work of the present paper is to realize exactly these pieces, the track, the locomotive, and the three
switches, as one tiling-independent cellular automaton.
